\subsection{Закон движения при РПД}
%%%%%%%%%%%%%%%%%%%%%%%%%%%%%%%%%%%%%%%%%%%%%%%%%%%%
%%%%%%%%%%%%%%%%%%%%%%%%%%%%%%%%%%%%%%%%%%%%%%%%%%%%
\z{Товарный поезд движется со скоростью $v_1  = 36 \, \frac{\text{км}}{\text{ч}}$. Спустя время $\tau = 30\, \text{мин}$ с той же станции в том же направлении вышел экспресс со скоростью $v_2  = 72\, \frac{\text{км}}{\text{ч}}$. Через время $t$ после выхода товарного поезда и на каком расстоянии $s$ от станции экспресс нагонит товарный поезд? Задачу решить, используя закон движения.}
%
{$t=1\,\text{ч};~s=36\, \text{км}$}

%%%%%%%%%%%%%%%%%%%%%%%%%%%%%%%%%%%%%%%%%%%%%%%%%%%%

 \z{Из пункта $A$ выехал велосипедист со скоростью $v_1  = 25\, \frac{\text{км}}{\text{ч}}$. Спустя время $t_0  = 6\, \text{мин}$ из пункта $B$, находящегося на расстоянии $L = 10\, \text{км}$ от пункта $A$, навстречу велосипедисту вышел пешеход. За время $t_2  = 50\, \text{с}$ пешеход прошел такой же путь, какой велосипедист проехал за $t_2  = 10\, \text{с}$. На каком расстоянии $s$ от пункта $A$ встретятся пешеход и велосипедист.%
}
%
{$s=8.75\, \text{км}$}
	
%%%%%%%%%%%%%%%%%%%%%%%%%%%%%%%%%%%%%%%%%%%%%%%%%%%%
\z{Из пунктов $A$ и $B$ одновременно навстречу друг другу выехали две машины. Через некоторое время они встретились и продолжили своё движение. Первая машина пришла в пункт назначения через $t_1  = 4\,\text{часа}$ после встречи, вторая --- через $t_2  = 1\, \text{час}$. Через какое время после выхода из пунктов $A$ и $B$ машины встретились?}
%
{$t=\sqrt{t_1\cdot t_2}=2\,\text{ч}$}

%%%%%%%%%%%%%%%%%%%%%%%%%%%%%%%%%%%%%%%%%%%%%%%%%%%%	
\z{На прямой дороге находятся велосипедист, мотоциклист и пешеход между ними. В начальный момент времени расстояние от пешехода до велосипедиста в 2 раза меньше, чем до мотоциклиста. Велосипедист и мотоциклист начинают двигаться навстречу друг другу со скоростями $20\, \frac{\text{км}}{\text{ч}}$ и $60 \, \frac{\text{км}}{\text{ч}}$ соответственно. В какую сторону и с какой скоростью должен идти пешеход, чтобы встретиться с велосипедистом и мотоциклистом в месте их встречи?%
}
%
{$v-6.7\, \frac{\text{км}}{\text{ч}}$, к велосипеду}

%%%%%%%%%%%%%%%%%%%%%%%%%%%%%%%%%%%%%%%%%%%%%%%%%%%%
\z{Экспериментатор Глюк наблюдал с безопасного расстояния за движением грозовой тучи. Увидев первую молнию, он засёк время и обнаружил, что услышал гром от неё только через $t_1  = 20\,\text{c}$. Через $\tau_1  = 3 \,\text{мин}$ после первой вспышки произошла вторая, а гром грянул с опозданием на $t_2  = 5\, \text{с}$. Подождав ещё $\tau_2  = 4\, \text{мин}$ после второй вспышки, Глюк увидел, как сверкнула последняя молния, и услышал звук грома от неё через $t_3  = 20\,\text{с}$. Определите скорость тучи $v$  и минимальное расстояние $h$ от Глюка за время наблюдения. Скорость звука в воздухе $u \approx 330\, \frac{\text{м}}{\text{с}}$, скорость света $c \approx 3\cdot10^8  \, \frac{\text{м}}{\text{с}}$.}
%
{$v=32\, \frac{\text{м}}{\text{с}};~h=1.4\,\text{км}$}

